% ----------------------------------------------------------
% Capítulo 2 — Referencial Teórico
% ----------------------------------------------------------
\chapter{Referencial Teórico}
\label{cap:referencial}

Este capítulo apresenta os catálogos de dados espaciais, o \textit{Tool Calling} e a Saída Estruturada, a execução local de LLMs e a orquestração com o LangChain, as métricas de avaliação e os trabalhos relacionados.

% ---
\section{Catálogos de Dados Espaciais}
% ---

A produção cartográfica brasileira é regida pelo Sistema Cartográfico Nacional (SCN), instituído pelo Decreto-Lei n.º 243/1967, que fixa as diretrizes e bases da cartografia brasileira e as escalas-padrão da cartografia sistemática terrestre \cite{scn2016}. O Decreto n.º 89.817/1984 estabeleceu as normas técnicas da cartografia nacional, inclusive os padrões de exatidão \cite{brasil1984decreto}. As folhas da cartografia sistemática são identificadas pelo Índice de Nomenclatura (INOM), um código alfanumérico como \mbox{SF-22-Y-D-II-4}, e pelo número do Mapa Índice (MI), como 2901 ou \mbox{2965-2-NE}. A Especificação Técnica para Produtos de Conjuntos de Dados Geoespaciais (ET-PCDG) \cite{etpcdg2016}, mantida pela DSG, detalha os requisitos técnicos dos produtos, como cartas topográficas, cartas ortoimagem, modelos digitais de elevação e ortoimagens.

A INDE é constituída por \textit{nós} provedores, pontos de acesso institucionais que catalogam e disponibilizam dados conforme o Perfil MGB \cite{perfilmgb2009}. A INDE e o BDGEx adotam os padrões do OGC para disponibilizar os dados, entre eles os serviços \textit{Web Map Service} (WMS), \textit{Web Feature Service} (WFS), \textit{Web Coverage Service} (WCS) e \textit{Catalogue Service for the Web} (CSW). O CSW \cite{ogccsw2007} é o padrão internacional para publicar e descobrir metadados de recursos geoespaciais. Suas operações de consulta, como \texttt{GetRecords}, buscam por título, data, tipo, área de abrangência geográfica e palavras-chave.

% ---
\section{\textit{Tool Calling}}
\label{sec:toolcalling}
% ---

O \textit{Tool Calling} (também chamado \textit{Function Calling}) é a capacidade de um LLM de reconhecer que a resposta a uma consulta exige uma função externa e de gerar a chamada dessa função, com os argumentos em JSON \cite{schick2023toolformer, qin2023toolllm}. O modelo passa a ser o componente que decide, num sistema que executa ferramentas determinísticas, como consultas a bancos de dados, chamadas a APIs e execução de código.

A Figura~\ref{fig:fluxo_toolcalling} mostra o fluxo. O sistema de orquestração (por exemplo, o LangChain) recebe a chamada gerada pelo modelo, executa a função e devolve o resultado ao modelo, que formula a resposta ao usuário. Este trabalho avalia a tradução, da consulta até a chamada da ferramenta com os seus argumentos.

\begin{figure}[h]
\centering
\begin{tikzpicture}[
  font=\footnotesize,
  node distance=0.9cm and 0.7cm,
  user/.style={rectangle, draw, rounded corners, minimum height=0.9cm, minimum width=2.2cm, align=center, fill=gray!8},
  llm/.style={rectangle, draw, rounded corners, minimum height=0.9cm, minimum width=2.2cm, align=center, fill=blue!10, very thick},
  payload/.style={rectangle, draw, rounded corners, minimum height=0.9cm, minimum width=2.2cm, align=center, fill=orange!12},
  tool/.style={rectangle, draw, rounded corners, minimum height=0.9cm, minimum width=2.2cm, align=center, fill=green!10},
  arr/.style={-{Stealth[length=2.2mm]}, thick}
]
  \node[user] (user1) {Usuário\\\textit{consulta}};
  \node[llm, right=of user1] (llm1) {LLM\\\textit{analisa}};
  \node[payload, right=of llm1] (call) {\textit{Tool Call}\\(JSON)};
  \node[tool, below=of call] (exec) {Ferramenta\\\textit{executa}};
  \node[payload, left=of exec] (result) {Resultado\\estruturado};
  \node[llm, left=of result] (llm2) {LLM\\\textit{responde}};
  \node[user, below=of llm2] (user2) {Usuário\\\textit{resposta}};

  \draw[arr] (user1) -- (llm1);
  \draw[arr] (llm1) -- (call);
  \draw[arr] (call) -- (exec);
  \draw[arr] (exec) -- (result);
  \draw[arr] (result) -- (llm2);
  \draw[arr] (llm2) -- (user2);
\end{tikzpicture}
\caption{Fluxo do mecanismo de \textit{Tool Calling}}
\label{fig:fluxo_toolcalling}
\fonte{Elaborado pelos autores.}
\end{figure}

No protocolo adotado pelo Ollama e pelo LangChain, cada ferramenta é definida por um \textit{schema} JSON com o nome da função, a sua descrição e os parâmetros aceitos, cada um com tipo e descrição. O modelo recebe essas definições junto com o \textit{prompt} de sistema, as instruções fixas que a aplicação lhe envia. A Figura~\ref{fig:exemplo_tool_schema} mostra uma definição simplificada da ferramenta de busca, com quatro dos doze parâmetros listados na Seção~\ref{sec:prototipo}.

\begin{figure}[h]
\centering
\begin{verbatim}
{
  "name": "buscar_catalogo",
  "description": "Busca produtos cartograficos no catalogo",
  "parameters": {
    "type": "object",
    "properties": {
      "keyword": {"type": "string",
                  "description": "Codigo MI, INOM ou nome da carta"},
      "scale":   {"type": "string",
                  "enum": ["1:25.000", "1:50.000", "1:100.000", ...],
                  "description": "Escala cartografica"},
      "state":   {"type": "string",
                  "description": "Nome do estado (normalizado)"},
      "productType": {"type": "string",
                       "description": "Tipo de produto SCN"}
    }
  }
}
\end{verbatim}
\caption{Exemplo simplificado de definição de ferramenta para busca em catálogo de produtos cartográficos}
\label{fig:exemplo_tool_schema}
\fonte{Elaborado pelos autores.}
\end{figure}

\subsection{\textit{Tool Calling} \textit{versus} Saída Estruturada}
\label{sec:tc-vs-se}

Na Saída Estruturada (\textit{Structured Output}), o modelo devolve a resposta num objeto JSON que segue um \textit{schema} predefinido, validado por bibliotecas como o Pydantic, em Python \cite{pydantic2019}, ou o Zod, em TypeScript. O modelo só decide quais campos preencher e com que valores. Bibliotecas como o Instructor \cite{instructor2023} conferem a saída contra o \textit{schema} e repetem a chamada quando ela é inválida.

A Saída Estruturada tem duas limitações estruturais. A primeira é que ela sempre devolve um objeto no formato do \textit{schema}. Mesmo com campos opcionais, toda consulta vira uma busca, e recusar exige um artifício no próprio \textit{schema}, como um campo de recusa. A segunda é que o modelo não consulta ferramentas durante a resposta. Uma verificação só lhe chega como nova tentativa, depois de emitido o objeto inteiro, como no Instructor. É esse o arranjo do controle da v3 (Seção~\ref{sec:v3}). Há ainda o risco de alucinações estruturadas, valores plausíveis sem apoio na consulta, inseridos para completar o \textit{schema}, hipótese que a Seção~\ref{sec:abordagens} testa.

No \textit{Tool Calling}, o modelo decide se e quando chama cada uma das ferramentas disponíveis. Ele pode responder sem chamar nenhuma ou encadear várias chamadas. Para este trabalho, o \textit{Tool Calling} pode ser mais adequado, porque deixa o modelo tratar consultas ambíguas, escolher os parâmetros relevantes e encadear validações antes da busca. Essa vantagem só se estabelece empiricamente. O Capítulo~\ref{cap:resultados} compara as duas abordagens numa só chamada (Seção~\ref{sec:abordagens}), na v3, em que as duas recebem o validador e só o \textit{Tool Calling} recebe as ferramentas auxiliares (Seção~\ref{sec:v3-resultados}), e no teste A/B, em acurácia e em tempo de processamento (Seção~\ref{sec:ab}).

% ---
\section{Modelos de Linguagem de Grande Porte e Execução Local}
% ---

Os LLMs, baseados na arquitetura \textit{Transformer} \cite{vaswani2017attention}, geram a resposta um \textit{token} (fragmento de palavra) por vez, condicionados ao \textit{prompt} de sistema, à consulta e ao que já geraram. O porte do modelo é o número de pesos, os parâmetros numéricos aprendidos no treinamento, contado em bilhões (B), como no Qwen 3 4B. A propriedade que mais importa aqui é seguir instruções em linguagem natural (\textit{instruction following}), obtida pelo ajuste fino com instruções e com \textit{feedback} humano \cite{ouyang2022training}. O \textit{Tool Calling} depende dela, porque o modelo precisa entender a consulta, interpretar o \textit{schema} das ferramentas e decidir como preenchê-lo.

\subsection{Vantagens e Desvantagens da Execução Local}

Executar os LLMs localmente, em vez de usar APIs de serviços em nuvem, traz vantagens para instituições como a DSG \cite{ollama2023}. As consultas e os dados não trafegam por redes externas, o que pesa em organizações com restrições de sigilo. Não há custo recorrente por \textit{token}. A organização mantém a versão exata do modelo em uso, o que permite auditá-lo. O sistema funciona sem conexão com a Internet.

A execução local tem também desvantagens:
\begin{itemize}
\item \textbf{Acesso limitado aos modelos de fronteira}. Os LLMs proprietários mais capazes não têm os pesos publicados e ocupam a maior parte das primeiras posições em \textit{benchmarks} de chamada de funções, como o \textit{Berkeley Function Calling Leaderboard} (BFCL) \cite{patil2025berkeley}.
\item \textbf{Custo de infraestrutura}. O gasto passa do pagamento por uso para o \textit{hardware}. A organização precisa adquirir e manter servidores com GPU e VRAM suficientes e arcar com energia, refrigeração e atualização de \textit{drivers}.
\item \textbf{Ciclo de atualização interno}. Cada nova versão de modelo exige repetir a avaliação sistemática.
\end{itemize}

\subsection{Ollama e Modelos com Suporte a \textit{Tool Calling}}

O Ollama \cite{ollama2023} é uma plataforma de código aberto que instala, executa e gerencia modelos de linguagem localmente. Ele expõe uma API compatível com o padrão OpenAI, o que facilita a integração com \textit{frameworks} como o LangChain. Para gerar as respostas, usa o \texttt{llama.cpp}, com os modelos no formato GGUF \cite{gguf2023}. O formato comporta quantizações como Q4\_0 e Q4\_K\_M, que reduzem a precisão dos pesos de 16 para cerca de 4 bits com perda pequena de qualidade, como mostram métodos de quantização pós-treinamento em 4 bits \cite{frantar2022gptq}. Este trabalho usou quantizações de 4 bits, Q4\_K\_M ou a equivalente da \textit{build} publicada de cada modelo (Seção~\ref{sec:selecao-modelos}), o que divide por cerca de quatro a memória de vídeo ocupada pelos pesos. Como ordem de grandeza, os pesos em 4 bits ocupam cerca de 0,5~GB por bilhão de parâmetros, e a memória do contexto, que cresce com o texto processado, se soma a eles.

A Tabela~\ref{tab:modelos_toolcalling} lista os quatro modelos avaliados (Seção~\ref{sec:selecao-modelos}). Os tamanhos orientaram a escolha de modelos executáveis na estação de referência, com GPU de 4~GB de VRAM, em que nenhum dos quatro coube inteiro. O excedente roda na CPU (\textit{offload}), o que aumenta a latência (Seção~\ref{sec:ambiente}).

\begin{table}[h]
\centering
\caption{Modelos avaliados, com suporte a \textit{Tool Calling} no Ollama}
\label{tab:modelos_toolcalling}
\resizebox{\textwidth}{!}{%
\begin{tabular}{|l|c|c|l|}
\hline
\textbf{Modelo} & \textbf{Parâmetros} & \textbf{Tamanho aprox. (4 bits)} & \textbf{Organização} \\
\hline
Gemma 4 E2B          & 2B efet.      & $\sim$4 GB$^{a}$ & Google DeepMind \\
Gemma 4 E4B          & 4B efet.      & $\sim$5{,}7 GB$^{a}$ & Google DeepMind \\
Qwen 3 4B            & 4B            & $\sim$2{,}3 GB & Alibaba Cloud \\
Mistral Nemo 12B     & 12B           & $\sim$6{,}6 GB  & Mistral AI e NVIDIA \\
\hline
\end{tabular}%
}
\fonte{Compilado pelos autores a partir de \citeauthoronline{ollama2023} (\citeyear{ollama2023}) e da documentação oficial dos modelos.}
\nota{$^{a}$~Tamanho da \textit{build} com quantização consciente de treinamento (QAT, 4 bits) publicada no Ollama (\texttt{gemma4:e2b-it-qat} e \texttt{gemma4:e4b-it-qat}). ``Efet.'' indica parâmetros efetivos; o total armazenado é maior (4,6B no E2B, segundo o Ollama), o que explica, em parte, o tamanho acima da ordem de grandeza de 0,5~GB por bilhão de parâmetros.}
\end{table}

% ---
\section{Orquestração com LangChain}
% ---

O LangChain \cite{langchain2022} é um \textit{framework} de código aberto para aplicações com LLMs, com uma interface comum para modelos, \textit{prompts} e ferramentas. Neste trabalho, ele liga a ferramenta de busca ao modelo, recolhe as chamadas de ferramenta que o modelo gera e é a interface única com os modelos no \textit{pipeline} de avaliação. Trocar a classe do Ollama (\texttt{ChatOllama}) pela de um provedor em nuvem (\texttt{ChatOpenAI} ou \texttt{ChatAnthropic}) exige alteração mínima no código, e a escolha entre execução local e em nuvem passa a ser uma decisão gerencial. O LangGraph \cite{langgraph2024}, extensão que modela fluxos como grafos de estados, não foi necessário. Na v1, a tradução é uma única decisão do modelo, e o laço da v3 é curto, tem limites fixos e é escrito explicitamente (Seções~\ref{sec:orquestracao} e~\ref{sec:v3}).

% ---
\section{Métricas de Avaliação}
\label{sec:metricas-conceitos}
% ---

As métricas servem para escolher o modelo e a abordagem a implantar, diante do compromisso entre qualidade e latência, e para reavaliar a solução quando surgirem modelos novos, o que o \textit{pipeline} de avaliação permite a custo baixo (Seção~\ref{sec:pipeline-avaliacao}). As métricas adotadas vêm da literatura de recuperação de informação \cite{manning2008introductionir}. Elas medem a correção total, a correção por campo e o tempo. A acurácia é o critério principal, e as demais desempatam ou medem o custo (Seção~\ref{sec:criterios-decisao}). A comparação com o gabarito é feita campo a campo, e não byte a byte, com as regras de normalização da Seção~\ref{sec:metricas}.

A \textbf{acurácia} é a proporção de consultas em que o modelo extraiu corretamente todos os parâmetros esperados, sem acrescentar parâmetros ausentes do gabarito. Nas consultas em que o gabarito prevê não buscar, a resposta é correta quando o modelo não chama a ferramenta de busca.
\[
\text{Acurácia} = \frac{\text{Número de consultas completamente corretas}}{\text{Total de consultas}}.
\]
É a métrica mais estrita, porque qualquer erro em qualquer parâmetro torna a consulta incorreta.

A análise por campo usa a \textbf{precisão}, o \textbf{\textit{recall}} e o \textbf{F1-\textit{score}},
\[
\text{Precisão} = \frac{\mathit{TP}}{\mathit{TP} + \mathit{FP}}
\qquad
\text{Recall} = \frac{\mathit{TP}}{\mathit{TP} + \mathit{FN}}
\qquad
F_1 = 2 \cdot \frac{\text{Precisão} \times \text{Recall}}{\text{Precisão} + \text{Recall}},
\]
em que $\mathit{TP}$ (\textit{true positives}) são os campos com valor aceito, $\mathit{FP}$ (\textit{false positives}) os campos com valor errado ou que não deviam ser emitidos e $\mathit{FN}$ (\textit{false negatives}) os campos obrigatórios omitidos. A decisão de recusar uma consulta é avaliada com as mesmas medidas, como uma classificação (Seção~\ref{sec:metricas}).

A \textbf{latência} é o tempo da tradução, do envio da consulta ao modelo até a decisão final de busca, numa ou em várias chamadas ao modelo. Ela é reportada por ambiente de execução, com estatísticas da sua distribuição sobre as consultas (Seção~\ref{sec:metricas}).

% ---
\section{Trabalhos Relacionados}
\label{sec:trabalhos-relacionados}
% ---

O ponto de partida direto deste trabalho, o protótipo do 1º CGEO \cite{cgeo2025prototipo}, é analisado na Seção~\ref{sec:prototipo}.

\paragraph{Interfaces de linguagem natural para bancos de dados.} Na formulação \textit{text-to-SQL}, o modelo gera a consulta SQL completa a partir da pergunta e do esquema do banco. O \textit{dataset} Spider \cite{yu2018spider} consolidou essa avaliação, com cerca de dez mil perguntas sobre 200 bancos. Com os LLMs, a abordagem passou a depender sobretudo do \textit{prompt}, e a ambiguidade e a subespecificação da linguagem natural estão entre os problemas centrais apontados pela literatura \cite{liu2025survey}. Neste trabalho, a formulação é mais restrita. O modelo não escreve SQL nem vê o esquema do banco. Ele só preenche os parâmetros tipados de uma ferramenta fixa, e o \textit{backend} monta com esses valores uma consulta SQL parametrizada. A restrição elimina o risco de SQL inválido ou injetado pelo modelo e permite validar cada argumento contra os enumerados do domínio. O custo é atender só às buscas que a ferramenta representa.

\paragraph{LLMs em geoinformação.} \citeonline{mai2023geospatialllm} discutem as oportunidades e os desafios dos modelos de fundação para a inteligência artificial geoespacial, e \citeonline{manvi2024geollm} mostram que o conhecimento geográfico dos LLMs serve a tarefas de predição quando o \textit{prompt} traz contexto do OpenStreetMap. Os LLMs passaram também a comandar tarefas de Sistemas de Informação Geográfica (SIG). O LLM-Geo, SIG autônomo construído por \citeonline{li2023autonomous} sobre o GPT-4, gera e executa o código que coleta, analisa e visualiza dados espaciais, e o GeoGPT \cite{zhang2024geogpt} planeja e encadeia ferramentas de SIG a partir de pedidos em linguagem natural. Nesses sistemas, o modelo tem alta autonomia. Neste trabalho, o modelo, aberto e executado localmente, só decide se busca e com quais argumentos, numa chamada na v1 e num laço curto com ferramentas auxiliares e um validador na v3, sempre com uma única ferramenta de busca.

\paragraph{Descoberta de dados geoespaciais.} A dificuldade de encontrar dados em geoportais é anterior aos LLMs. \citeonline{hu2015metadata} a atribuem, em parte, à heterogeneidade dos metadados e propõem harmonização de tópicos e busca semântica em geoportais baseados em \textit{Linked Data}. Com os LLMs, \citeonline{ning2025autonomous} propõem um agente que escolhe a fonte de dados numa lista predefinida e gera o programa que a recupera, e \citeonline{li2025question} combinam LLMs e um grafo de conhecimento construído a partir de cinco grandes portais para responder a perguntas de descoberta de dados. No caso da DSG, o catálogo é único e já estruturado, e o problema se reduz a traduzir a intenção do usuário nos filtros desse catálogo. Por isso a solução dispensa grafo de conhecimento e geração de código.

\paragraph{Avaliação da chamada de ferramentas.} Os trabalhos sobre uso de ferramentas por LLMs trouxeram \textit{benchmarks} próprios, como o API-Bank \cite{li2023apibank}, o ToolBench do ToolLLM \cite{qin2023toolllm}, com mais de 16 mil APIs reais, e o APIBench do Gorilla \cite{patil2023gorilla}. O BFCL \cite{patil2025berkeley} compara a chamada gerada com a esperada pela árvore sintática e avalia também a abstenção. Segundo ele, os melhores modelos vão bem em chamadas de um único turno, mas ainda não em memória, decisões dinâmicas e raciocínio de longo prazo. Esses \textit{benchmarks}, na maioria em inglês, de domínio geral e com muitas ferramentas, não cobrem o cenário deste trabalho, que reúne uma única ferramenta de busca, consultas em português com o vocabulário da cartografia sistemática (MI, INOM, escalas, centros produtores) e modelos pequenos executados localmente. Por isso este trabalho construiu um \textit{dataset} próprio, com gabarito auditado e comparação campo a campo semelhante, em espírito, à comparação estrutural do BFCL.

\paragraph{Ferramentas auxiliares e retorno.} No ReAct \cite{yao2023react}, o modelo alterna raciocínio e ações, e o resultado de cada ação alimenta o passo seguinte. A correção funciona quando o sinal vem de fora do modelo. O CRITIC \cite{gou2024critic} verifica a saída com ferramentas e a revisa, e o \textit{self-debugging} \cite{chen2024selfdebug} corrige consultas a partir do resultado da execução. A autocorreção sem sinal externo \cite{madaan2023selfrefine} pode não ajudar e até piorar o raciocínio \cite{huang2024cannot}. O ToolSandbox \cite{lu2025toolsandbox} aponta a normalização de argumentos, como datas relativas e lugares, entre as maiores dificuldades. O When2Call \cite{ross2025when2call} mostra que decidir entre chamar, perguntar ou recusar é uma capacidade própria, fraca em modelos pequenos. O $\tau$-bench \cite{yao2025taubench} avalia diálogos com usuário simulado, e o CRITICTOOL \cite{huang2025critictool} classifica os erros de chamada que o modelo precisa perceber e corrigir.

Ajustar um sistema pelos erros de um conjunto rotulado arrisca torná-lo bom só nesse conjunto (sobreajuste). O risco vale para as lições que o modelo acumula a partir de sucessos e falhas, como no Reflexion \cite{shinn2023reflexion} e no ExpeL \cite{zhao2024expel}, e para a otimização automática de \textit{prompts} \cite{yang2024opro,khattab2024dspy}. Medir o resultado exige um conjunto de teste que o desenvolvimento não tenha visto \cite{dwork2015reusable}. A v3 (Seção~\ref{sec:v3}) leva a um domínio restrito as ferramentas auxiliares determinísticas, a validação com retorno e o pedido de esclarecimento a um usuário simulado, acrescentados em degraus, e foi testada num lote que o desenvolvimento não leu.
